\documentclass[12pt]{article}
\usepackage[utf8]{inputenc}
\usepackage{amsmath}
\usepackage{amsfonts}
\usepackage{amssymb}
\usepackage{geometry}
\usepackage{booktabs}
\usepackage{array}
\usepackage{float}

\geometry{margin=1in}

\title{Insurance Costs Methodology\\Comprehensive Transmission Cost Calculator}
\author{Andrew Igdal}
\date{\today}

\begin{document}

\maketitle

\section{Overview}
This document describes the methodology for calculating insurance costs for transmission line projects. The calculation considers both AC and DC transmission projects, accounting for different infrastructure components and their associated insurance premiums.

\section{Methodology}

\subsection{Project Types}
The insurance cost calculation accommodates two types of transmission projects:
\begin{itemize}
    \item \textbf{AC Projects}: Traditional alternating current transmission lines
    \item \textbf{DC Projects}: Direct current transmission lines with converter stations
\end{itemize}

\subsection{Infrastructure Components}
The calculation considers the following infrastructure components:
\begin{enumerate}
    \item \textbf{Conductor}: The transmission line conductors
    \item \textbf{Structure}: Transmission towers and support structures
    \item \textbf{Converter}: DC converter stations (only for DC projects)
\end{enumerate}

\subsection{Mathematical Formulation}

The annual insurance payment per mile of transmission line is calculated as:
\begin{equation}
P_{mile} = (C_{conductor} + C_{structure}) \times \pi_{insurance}
\end{equation}

The total annual line payment is:
\begin{equation}
P_{line} = P_{mile} \times L
\end{equation}

The annual converter payment (for DC projects only) is:
\begin{equation}
P_{converter} = C_{converter} \times \pi_{insurance} \times \phi_{converter}
\end{equation}

The total annual insurance payment is:
\begin{equation}
P_{annual} = P_{line} + P_{converter}
\end{equation}

The total project lifetime insurance payment is:
\begin{equation}
P_{lifetime} = P_{annual} \times T
\end{equation}

The present value of total project lifetime insurance payments is:
\begin{equation}
PV_{insurance} = \sum_{t=1}^{T} \frac{P_{annual}}{(1 + r)^t}
\end{equation}

Where:
\begin{itemize}
    \item $C_{conductor}$ = Conductor cost (USD)
    \item $C_{structure}$ = Structure cost (USD)
    \item $C_{converter}$ = Converter cost (USD)
    \item $\pi_{insurance}$ = Insurance premium rate (decimal)
    \item $\phi_{converter}$ = Binary variable (0 for AC, 1 for DC)
    \item $L$ = Project length in miles
    \item $T$ = Project lifetime in years
    \item $r$ = Firm interest rate (decimal)
    \item $t$ = Time period (year)
\end{itemize}

\section{Input Parameters}

\begin{table}[H]
\centering
\caption{Required Input Parameters}
\begin{tabular}{@{}ll@{}}
\toprule
\textbf{Parameter} & \textbf{Description} \\
\midrule
Conductor Cost & Cost of transmission line conductors (USD) \\
Structure Cost & Cost of transmission towers and structures (USD) \\
Converter Cost & Cost of DC converter stations (USD) \\
Phi Converter & Binary variable: 0 = AC project, 1 = DC project \\
Firm Interest Rate & Discount rate for present value calculation (decimal) \\
Risk-Free Interest Rate & Risk-free interest rate (decimal) \\
Project Lifetime & Expected project lifetime in years \\
Insurance Premium & Annual insurance premium rate (decimal) \\
Project Length & Total length of transmission line in miles \\
\bottomrule
\end{tabular}
\end{table}

\section{Outputs}

The script provides the following outputs:
\begin{itemize}
    \item Annual insurance cost
    \item Total project lifetime insurance cost
    \item Present value of total project lifetime insurance cost
\end{itemize}

\section{Assumptions}

\begin{itemize}
    \item Insurance premiums are paid annually
    \item The insurance premium rate remains constant over the project lifetime
    \item The firm interest rate remains constant over the project lifetime
    \item All costs are expressed in nominal USD
    \item Converter costs are only applicable to DC projects
\end{itemize}

\section{Usage Notes}

\begin{itemize}
    \item Input all cost values in USD
    \item Interest rates should be entered as decimals (e.g., 0.05 for 5\%)
    \item Insurance premium should be entered as a decimal (e.g., 0.02 for 2\%)
    \item Phi converter should be entered as 0 for AC projects or 1 for DC projects
    \item Project length should be entered in miles
    \item Project lifetime should be entered in years
\end{itemize}

\section{Example Calculation}

For a 100-mile DC transmission project with:
\begin{itemize}
    \item Conductor cost: \$50,000/mile
    \item Structure cost: \$30,000/mile
    \item Converter cost: \$500,000
    \item Insurance premium: 2\%
    \item Project lifetime: 30 years
    \item Firm interest rate: 5\%
\end{itemize}

The annual insurance cost would be:
\begin{align}
P_{mile} &= (50,000 + 30,000) \times 0.02 = \$1,600/\text{mile} \\
P_{line} &= 1,600 \times 100 = \$160,000 \\
P_{converter} &= 500,000 \times 0.02 \times 1 = \$10,000 \\
P_{annual} &= 160,000 + 10,000 = \$170,000
\end{align}

\end{document}
